\section*{Notation}
\addcontentsline{toc}{section}{Notation}
\label{sec:notation}

The symbols below are grouped by the pipeline stage in which they are
introduced and are used consistently throughout the document. Terms rather than
symbols (SUMO, TAZ, OD matrix, deviation plan, baseline, MCDA) are defined in
Table~\ref{tab:method_concepts}. Where a symbol carries a run index it is
written as a parenthesised superscript, $x^{(r)}$, and where it refers to a
candidate plan it carries the subscript $k$, with $k=0$ reserved for the
baseline condition.

\vspace{0.6em}
\noindent\textbf{Network and graph}
\nopagebreak

\begin{longtable}{@{}p{0.16\linewidth}>{\raggedright\arraybackslash}p{0.60\linewidth}p{0.16\linewidth}@{}}
\toprule
Symbol & Meaning & Introduced \\
\midrule
\endfirsthead
\toprule
Symbol & Meaning & Introduced \\
\midrule
\endhead
\bottomrule
\endfoot
\bottomrule
\endlastfoot

$G=(V,E,\ell)$ & Weighted directed junction-level graph of the road network \newline & \S\ref{subsec:deviation_plan_maker} \\
$V$ & Set of junctions (graph nodes) & \S\ref{subsec:deviation_plan_maker} \\
$E$ & Set of road edges (graph arcs) & \S\ref{subsec:deviation_plan_maker} \\
$\ell(e)$ & Length of edge $e$ in meters; $\ell(P)$ is the total length of path $P$ & \S\ref{subsec:deviation_plan_maker} \\
$n_L(e)$ & Number of lanes of edge $e$ & Eq.~\eqref{eq:zone_mass} \\
$\deg(v)$ & Number of edges incident to junction $v$ in the queried direction & \S\ref{sssec:junction_graph} \\
$G'$ & Residual graph $(V, E\setminus B)$ used for detour search & Eq.~\eqref{eq:deviation_plan} \\
\addlinespace
\multicolumn{3}{@{}l}{\textbf{Demand generation}}\\
\addlinespace
$i,\,j$ & Origin and destination zone indices & \S\ref{ssec:stage_1_demand_generation} \\
$\mathcal{E}_i$ & Set of road edges contained in zone $i$ & Eq.~\eqref{eq:zone_mass} \\
$m_i$ & Zone mass of zone $i$, a lane-kilometre proxy for zonal activity & Eq.~\eqref{eq:zone_mass} \\
$N$ & Total daily demand, in trips & \S\ref{par:temporal_demand_profile} \\
$h$ & Hour of the simulated day, $h=0,\ldots,23$ & \S\ref{par:temporal_demand_profile} \\
$p_h$ & Share of the daily demand departing in hour $h$ & Eq.~\eqref{eq:twopeak_gaussian} \\
$n_h$ & Integer number of trips departing in hour $h$; $\sum_h n_h = N$ & \S\ref{par:temporal_demand_profile} \\
$g(h;\mu,\sigma)$ & Gaussian kernel centred on hour $\mu$ with spread $\sigma$ & Eq.~\eqref{eq:twopeak_gaussian} \\
$\mu_{AM},\,\mu_{PM}$ & Morning and evening peak hours & Eq.~\eqref{eq:twopeak_gaussian} \\
$a_{AM},\,a_{PM}$ & Amplitudes of the morning and evening peaks & Eq.~\eqref{eq:twopeak_gaussian} \\
$\sigma$ & Temporal spread of the demand peaks, in hours & Eq.~\eqref{eq:twopeak_gaussian} \\
$c_{ij}$ & Travel-cost impedance from zone $i$ to zone $j$, in meters & Eq.~\eqref{eq:gravity_model} \\
$\beta$ & Distance-decay parameter of the gravity model & Eq.~\eqref{eq:gravity_model} \\
$\gamma_{ij}$ & Unnormalised gravity weight of OD pair $(i,j)$ & Eq.~\eqref{eq:gravity_model} \\
$\Theta_{ij}$ & Normalised OD share matrix; $\sum_{ij}\Theta_{ij}=1$ & Eq.~\eqref{eq:od_share} \\
$O_i$ & Origin (production) share of zone $i$ & Eq.~\eqref{eq:od_share} \\
$T^h_{ij}$ & Integer number of trips from zone $i$ to zone $j$ in hour $h$ & \S\ref{par:gravity_od_model} \\
$\rho_i,\,\varepsilon_i$ & Residential and employment weights of zone $i$ (commute overlay) & \S\ref{par:commute_overlay} \\
$\omega$ & Commute-overlay mixing strength (\texttt{peak\_weight}); $\omega=0$ recovers the base model & Eq.~\eqref{eq:commute_blend} \\
\addlinespace
\multicolumn{3}{@{}l}{\textbf{Route library and trip instantiation}}\\
\addlinespace
$k$ & Number of shortest paths requested from Yen's algorithm & \S\ref{par:k_shortest_candidate_routes} \\
$n_s$ & Number of sampled (origin-edge, destination-edge) pairs per OD pair & \S\ref{par:k_shortest_candidate_routes} \\
$\mathcal{L}_{ij}$ & Route library: candidate paths retained for OD pair $(i,j)$ & \S\ref{par:k_shortest_candidate_routes} \\
$\pi$ & A path, as an ordered edge sequence & \S\ref{par:route_selection} \\
$t^{\mathrm{dep}}_i$ & Departure time of trip $i$, drawn uniformly within its hour & \S\ref{par:departure_time} \\
\addlinespace
\multicolumn{3}{@{}l}{\textbf{Closure and deviation plans}}\\
\addlinespace
$E_c$ & Ordered set of closed edges $(e^c_1,\ldots,e^c_n)$ & \S\ref{sssec:detour_endpoints} \\
$B$ & Blocked set, $B \supseteq E_c$: edges the detour search may not use & Eq.~\eqref{eq:deviation_plan} \\
$s$ & Detour origin: junction at which a plan begins & Eq.~\eqref{eq:deviation_plan} \\
$t$ & Detour target: junction at which a plan rejoins the network & Eq.~\eqref{eq:deviation_plan} \\
$s_0,\,t_0$ & Natural (pre-promotion) endpoints of the closure & \S\ref{sssec:detour_endpoints} \\
$P$ & Detour path $(a_1,\ldots,a_{|P|})$ from $s$ to $t$ in $G'$ & Eq.~\eqref{eq:deviation_plan} \\
$d$ & BFS depth of a plan origin relative to the effective detour source & Eq.~\eqref{eq:deviation_plan} \\
$d_{\max}$ & Maximum upstream search depth & \S\ref{par:upstream} \\
$\mathcal{D}$ & A deviation plan, the tuple $(s,t,P,\ell(P),d)$ & Eq.~\eqref{eq:deviation_plan} \\
$\mathcal{D}_k$ & The $k$-th candidate deviation plan & \S\ref{sssec:mc_design} \\
$\mathcal{D}^*$ & Plan selected for a given vehicle during route rewriting & \S\ref{sssec:route_rewriting} \\
$\mathcal{P}$ & Ordered list of paths returned by a candidate-generation method & \S\ref{par:yen_dpm} \\
$P_i$ & The $i$-th shortest path accepted by Yen's algorithm & \S\ref{par:yen_dpm} \\
$u$ & An upstream junction discovered by the reverse BFS, at depth $d(u)$ & \S\ref{par:upstream} \\
$K$ & Number of candidate plans evaluated & \S\ref{sssec:mc_iteration} \\
\addlinespace
\multicolumn{3}{@{}l}{\textbf{Monte Carlo simulation}}\\
\addlinespace
$r$ & Monte Carlo run index (also the random seed of that run) & \S\ref{sssec:mc_design} \\
$R$ & Number of Monte Carlo runs retained for the paired analysis & \S\ref{sssec:mc_convergence} \\
$\mathcal{T}$ & Fixed trip file (\texttt{trips.xml}) shared by every run & Eq.~\eqref{eq:duarouter_run} \\
$\alpha$ & Upper bound of \texttt{duarouter}'s randomised edge-cost multiplier & Eq.~\eqref{eq:duarouter_run} \\
$\mathcal{R}^{(r)}_0$ & Baseline route realisation of run $r$ & Eq.~\eqref{eq:duarouter_run} \\
$\mathcal{R}^{(r)}_k$ & Route realisation of run $r$ after applying plan $\mathcal{D}_k$ & \S\ref{sssec:mc_iteration} \\
$\mathcal{Q}_{k,v}$ & Edge sequence driven by vehicle $v$ under plan $\mathcal{D}_k$ & Eq.~\eqref{eq:evaluator_impacted_set} \\
$\mathcal{M}$ & Edge-endpoint map used when splicing a detour into a route & \S\ref{sssec:mc_iteration} \\
$M^{(r)}_k$ & Metric vector of condition $k$ in run $r$ & \S\ref{sssec:mc_metrics} \\
$\nu^{(r)}_k$ & Number of vehicles rerouted by plan $\mathcal{D}_k$ in run $r$ & \S\ref{sssec:mc_iteration} \\
$[T_{\mathrm{start}}, T_{\mathrm{end}}]$ & Simulated time window & \S\ref{sssec:mc_iteration} \\
$\mathcal{V}_{\mathrm{done}}$ & Set of vehicles completing their trip within the window & Eq.~\eqref{eq:mc_delay} \\
$\theta_v$ & SUMO \texttt{timeLoss} of vehicle $v$, in seconds & Eq.~\eqref{eq:mc_delay} \\
$\tau_v$ & Door-to-door duration of trip $v$, in seconds & Eq.~\eqref{eq:mc_tt} \\
$D$ & Total time-loss delay over completed trips, in vehicle-hours & Eq.~\eqref{eq:mc_delay} \\
$\bar{\tau}$ & Mean travel time per completed trip, in minutes & Eq.~\eqref{eq:mc_tt} \\
$C^{(r)}_k$ & Number of trips completed under condition $k$ in run $r$ & Eq.~\eqref{eq:evaluator_feasibility} \\
\addlinespace
\multicolumn{3}{@{}l}{\textbf{Evaluation and decision support}}\\
\addlinespace
$\Delta D^{(r)}_k$ & Paired time-loss delay difference, $D^{(r)}_k - D^{(r)}_0$ & Eq.~\eqref{eq:paired_delta} \\
$\overline{\Delta D_k}$ & Mean paired delay difference over the $R$ runs & Eq.~\eqref{eq:paired_delta_mean} \\
$\Delta\bar{\tau}^{(r)}_k$ & Paired mean-travel-time difference & Eq.~\eqref{eq:evaluator_delta_travel_time} \\
$\bar{x},\,s_x,\,\tilde{x}$ & Sample mean, sample standard deviation and median of a paired-difference vector & Eq.~\eqref{eq:evaluator_descriptive_stats} \\
$t_{0.975,\,R-1}$ & Student-$t$ critical value for a two-sided 95\% interval & Eq.~\eqref{eq:evaluator_t_interval} \\
$\mathcal{I}_k$ & Set of vehicles impacted by plan $\mathcal{D}_k$ & Eq.~\eqref{eq:evaluator_impacted_set} \\
$\Delta T^{(r)}_{k,\mathcal{I}}$ & Total additional travel time of impacted users, in vehicle-minutes & Eq.~\eqref{eq:evaluator_impacted_delta_time} \\
$\phi_{\mathrm{net}}$ & Fraction of impacted users whose travel time increased & Eq.~\eqref{eq:evaluator_phi} \\
$q$ & Decision-criterion index & Eq.~\eqref{eq:evaluator_normalisation} \\
$x_{kq}$ & Raw value of plan $k$ on criterion $q$ & Eq.~\eqref{eq:evaluator_normalisation} \\
$z_{kq}$ & Min-max normalised value of plan $k$ on criterion $q$ & Eq.~\eqref{eq:evaluator_normalisation} \\
$w_q$ & Weight of criterion $q$; $w_q \geq 0$ and $\sum_q w_q = 1$ & Eq.~\eqref{eq:evaluator_additive_score} \\
$S_k$ & Additive weighted score of plan $k$; lower is better & Eq.~\eqref{eq:evaluator_additive_score} \\
$N_w$ & Number of weight vectors sampled in the sensitivity analysis & \S\ref{sssec:evaluator_mcda} \\

\end{longtable}

\noindent Two conventions are worth stating explicitly. First, all three
decision criteria of Equation~\eqref{eq:evaluator_criteria_vector} are
expressed as \emph{costs}, so lower values are preferred everywhere and a
positive paired difference always means that a plan worsened the metric.
Second, $k$ denotes the number of shortest paths requested from Yen's
algorithm in Sections~\ref{par:k_shortest_candidate_routes}
and~\ref{par:yen_dpm}, and the plan index thereafter; the two coincide for
Yen-generated families, where the $k$ enumerated paths become the $K$ evaluated
plans.
